% =====================================================================
%  3 Zeitplanung  (Aufgabenstellung 1.1: „hinreichend detaillierter
%  Zeitplan und Erfassung wesentlicher Arbeitspakete und Meilensteine“)
%  ACHTUNG Bewertung: „Wurde der initial festgelegte Zeitplan eingehalten?“
%  -> realistisch planen; Abweichungen später im Abstract reflektieren.
%  Feedback der Tutor:innen kann bis zu 7 Tage dauern -> einplanen!
% =====================================================================
\section{Zeitplanung}
\TODO{Kurze Hinführung; Planungszeitraum; verfügbare Wochenstunden}

\subsection{Arbeitspakete}
\begin{longtable}{@{}L{1.2cm}L{6.3cm}L{1.6cm}L{1.4cm}L{4.6cm}@{}}
  \caption{Arbeitspakete}\label{tab:arbeitspakete}\\
  \toprule
  \textbf{AP} & \textbf{Arbeitspaket} & \textbf{Phase} & \textbf{Std.} & \textbf{Ergebnis} \\
  \midrule
  \endfirsthead
  \toprule
  \textbf{AP} & \textbf{Arbeitspaket} & \textbf{Phase} & \textbf{Std.} & \textbf{Ergebnis} \\
  \midrule
  \endhead
  \bottomrule
  \endfoot
  AP1  & Projektplanung & 1 & \TODO{} & Projektdokumentation \\
  AP2  & Anforderungsermittlung & 1 & \TODO{} & Anforderungsdokument \\
  AP3  & Spezifikation & 1 & \TODO{} & Spezifikationsdokument \\
  AP4  & Feedback Phase 1 einarbeiten & 2 & \TODO{} & überarbeitete Dokumente \\
  AP5  & Architektur und Design & 2 & \TODO{} & Architekturdokument \\
  AP6  & Implementierung & 2 & \TODO{} & Quellcode (GitHub) \\
  AP7  & Feedback Phase 2 einarbeiten & 3 & \TODO{} & überarbeitete Dokumente \\
  AP8  & Testen und Qualitätssicherung & 3 & \TODO{} & Testdokument \\
  AP9  & Benutzeranleitung und Auslieferung & 3 & \TODO{} & ausführbare Anwendung \\
  AP10 & Projektabschluss und Reflexion & 3 & \TODO{} & Abstract, finale Abgabe \\
\end{longtable}
\quelle{Eigene Darstellung.}

\subsection{Meilensteine}
\begin{longtable}{@{}L{1.4cm}L{10.0cm}L{3.8cm}@{}}
  \caption{Meilensteine}\label{tab:meilensteine}\\
  \toprule
  \textbf{ID} & \textbf{Meilenstein} & \textbf{Geplant bis} \\
  \midrule
  \endfirsthead
  \bottomrule
  \endfoot
  M1 & Abgabe Portfolioteil 1 (Konzeptionsphase) & \TODO{Datum} \\
  M2 & Abgabe Portfolioteil 2 (Erarbeitungs- und Reflexionsphase) & \TODO{Datum} \\
  M3 & Funktionsfähige Anwendung, Tests abgeschlossen & \TODO{Datum} \\
  M4 & Abgabe Portfolioteil 3 (Finalisierungsphase) & \TODO{Datum} \\
\end{longtable}
\quelle{Eigene Darstellung.}

\subsection{Zeitplan}
\TODO{Gantt-Diagramm der Arbeitspakete (z. B. mit draw.io) als docs/diagramme/export/zeitplan.pdf}
\begin{figure}[H]
  \caption{Zeitplan der Arbeitspakete}\label{abb:zeitplan}
  \centering
  \grafik{zeitplan}
  \quelle{Eigene Darstellung.}
\end{figure}
